% ======================================================================
% 3.6 Настройка в реалистичных условиях. Ответственный: C — Дучин Кирилл.
%
% Что сдать (main.tex: теория 1702–1734, порядок 1736–1760, пример 1762–1831;
% требования к отчёту 1880–1884):
%   1. Явно задать: функцию трафика lambda(t) (суточная синусоида + всплеск),
%      значение задержки tau, уравнение регулятора с задержкой, начальную
%      историю на [-tau, 0], T_sim, Delta t и способ учёта задержки
%      (целое m=tau/Delta t или интерполяция).
%   2. Вычислить средние затраты за сутки модельного времени (формула
%      трапеций по L_n = a*c_n + b*q_n).
%   3. Найти минимум средних затрат градиентным спуском с численным градиентом,
%      начав с u0* и k* из п. 3.3–3.4; линии уровня на плоскости (u0,k)
%      и траектория спуска.
%   4. Сравнить найденные настройки с аналитическими и ответить на сквозной
%      вопрос работы; объяснить ограничения выводов.
%
% Параметры варианта 3 (приложение Б): lambda=600, mu=80, a=1, b=1.
% ВАЖНО (main.tex, строки 1825–1831): u0* — лишь исходная настройка, k* — порог
% колебаний модели без задержки, а не минимум затрат; выпуклость целевой функции
% не установлена; шаг Delta t, пригодный без задержки, не переносится автоматически
% на систему с задержкой — устойчивость и точность проверить отдельно.
% Во всех прогонах для численного градиента — одинаковые трафик, история, сетка.
%
% Запуск: python experiments.py 3.6.
% Рисунки: figures/3-6-levels.pdf, figures/3-6-dynamics.pdf.
% ======================================================================

\subsection{Настройка в реалистичных условиях}

\taskhead{Постановка эксперимента}
\todo{C: явно выписать $\lambda(t)$ (амплитуда суточного колебания, всплеск: интервал и величина), значение $\tau$, уравнение регулятора с задержкой, начальную историю, $T_{\mathrm{sim}}$, $\Delta t$ и проверку его устойчивости/точности; способ учёта задержки.}

\taskhead{Средние затраты и численный градиент}
\todo{C: описать вычисление $\overline\Phi(u_0,k)$ (формула трапеций по $L_n=ac_n+bq_n$); численный градиент отдельными прогонами с одинаковыми трафиком/историей/сеткой; шаги $\varepsilon_u,\varepsilon_k$; критерий остановки; допустимость $0<u_0<1$, $k>0$; что делать с неустойчивой траекторией.}

\taskhead{Результаты оптимизации}
\todo{C: листинг вызовов и запуск python experiments.py 3.6; старт с $u_0^*,k_*$ из п. 3.3–3.4.}
\begin{table}[htbp]
\centering\small
\caption{Исходные и найденные настройки регулятора}
\label{tab:3-6}
\begin{tabular}{|l|r|r|r|}\hline
Настройки & $u_0$ & $k$ & $\overline\Phi$\\\hline
Исходные (из п. 3.3–3.4) & — & — & —\\
Найденные спуском & — & — & —\\
\hline\end{tabular}
\end{table}
\taskfigure[0.8]{figures/3-6-levels.pdf}{Линии уровня средних затрат в плоскости $(u_0,k)$ и траектория спуска}{fig:3-6-levels}
\taskfigure[0.85]{figures/3-6-dynamics.pdf}{Трафик $\lambda(t)$, очередь $q(t)$ и ресурс $c(t)$ за сутки модельного времени}{fig:3-6-dynamics}
\todo{C: пояснить графики; указать параметры численного градиента и критерий остановки.}

\taskhead{Вывод}
\todo{C: сравнить найденные настройки с аналитическими; ответить на сквозной вопрос работы (как настраивать регулятор, чтобы минимизировать затраты при сохранении устойчивости); ограничения выводов: локальный характер минимума, зависимость от выбранного трафика и $\tau$.}
